% Generated from locale/tr/content/first-order-logic/introduction/semantic-notions.tex; do not hand-edit.


\olfileid[tr]{fol}{int}{sem}

\olsection{Semantik Kavramlar}

Yukarıda $\Sat M{!A}[s]$ olup olmadığını ele alırken, kolaylık olsun
diye~$s$'nin bütün !!{variable}s ifadelerine değerler atamasına izin
verdiğimizi, fakat yalnızca~$!A$ içindeki !!{variable}s ifadelerine
atanan değerlerin kullanıldığını söyledik. Aslında yalnızca~$!A$
içindeki \emph{serbest} değişkenlerin değerleri önemlidir. Titiz
davrandığımız için elbette bu olguyu ispatlayacağız. !!{sentence}s
serbest değişken içermediğinden, bir yapı içinde sağlanıp
sağlanmamaları bakımından~$s$ hiçbir önem taşımaz. Dolayısıyla $!A$
!!a{sentence} olduğunda $\Sat M{!A}$ ifadesini ``bütün~$s$ atamaları
için $\Sat M{!A}[s]$'' anlamında tanımlayabiliriz. Bu da en az bir~$s$
için $\Sat M{!A}[s]$ olmasıyla denktir. !!{formula}s için çalışan bir
sağlama tanımı elde etmek üzere !!{variable} atamalarını tanıtmamız
gerekir; fakat !!{sentence}s için sağlama, !!{variable} atamalarından
bağımsızdır.

``$\Sat M{!A}$'' ifadesini tanımladıktan sonra birinci dereceden
mantığın !!{sentence}s ifadeleri açısından ``durum'' ve ``içinde
doğru'' kavramlarının ne anlama geldiğini biliriz. Ardından
!!{sentence}s için $\Sat M{!A}$ tanımına dayanarak temel semantik
geçerlilik, mantıksal gerektirme ve sağlanabilirlik kavramlarını
tanımlayabiliriz. Bir tümce, her !!{structure} onu sağlıyorsa geçerlidir:
$\Entails !A$. Bir~$!A$ tümcesi,~$\Gamma$ içindeki bütün !!{sentence}s
ifadelerini sağlayan her !!{structure} aynı zamanda~$!A$'yı da
sağlıyorsa bir !!{sentence}s kümesi~$\Gamma$ tarafından mantıksal
olarak gerektirilir: $\Gamma\Entails !A$. Bir !!{sentence}s kümesi de,
bir !!{structure} içindeki bütün !!{sentence}s ifadelerini aynı anda
sağlıyorsa sağlanabilirdir.

!!{formula}s tümevarımlı olarak tanımlanmış ve sağlama da tek tek
!!{formula}s ele alındığında bunların yapısı üzerinde tümevarımla
tanımlandığından, semantiğimizin
özelliklerini ispatlamak ve tanımlanan semantik kavramları birbiriyle
ilişkilendirmek için tümevarımı kullanabiliriz. Bu özelliklerden
bazılarını derleyip ispatlayacağız. Bunu kısmen her biri kendi başına
ilginç olduğu için, fakat esas olarak semantikle !!{derivation}
sistemleri arasındaki ilişkiyi incelemeye geçtiğimizde birçoğu yararlı
olacağı için yapacağız. Bunun için henüz söz etmediğimiz bazı sözdizimsel
kavram ve işlemleri de (kesin biçimde, yani tümevarımla) tanımlamamız
gerekecektir.

